\documentclass[11pt]{article}

\usepackage[a4paper,margin=1in]{geometry}
\usepackage{booktabs}
\usepackage{array}
\usepackage{longtable}
\usepackage{enumitem}
\usepackage{xcolor}
\usepackage{hyperref}
\usepackage{fancyhdr}
\usepackage{titlesec}
\usepackage{caption}

\hypersetup{
    colorlinks=true,
    linkcolor=blue,
    urlcolor=blue,
    citecolor=blue
}

\pagestyle{fancy}
\fancyhf{}
\fancyhead[L]{Final Report}
\fancyhead[R]{CSE406 Security Sessional}
\fancyfoot[C]{\thepage}

\newcommand{\worked}{\textcolor{green!50!black}{\textbf{As expected}}}
\newcommand{\failed}{\textcolor{red!70!black}{\textbf{Differed from expectation}}}
\newcommand{\partialres}{\textcolor{orange!80!black}{\textbf{Partially / conditionally}}}
\newcommand{\open}{\textcolor{gray}{\textbf{Unresolved / under investigation}}}

\title{
    \textbf{ICMP Blind Connection-Reset \& Blind Throughput Reduction\\
    Attack Against TCP} \\
    \vspace{0.3cm}
    \large Final Implementation \& Evaluation Report \\
    \vspace{0.2cm}
    \large CSE 406 --- Computer Security Sessional
}
\author{
    Anika Morshed (2105068) \\
    Diganta Saha Tirtha (2105081) \\
    Section: B1
}
\date{\today}

\begin{document}

\maketitle

\begin{abstract}
This report documents the implementation and empirical evaluation of two
ICMP-based blind injection attacks against TCP --- a connection-reset
attack (Type 3, Code 3) and a Path-MTU throughput-reduction attack
(Type 3, Code 4) --- together with two candidate defenses against them.
For each attack and each defense we record the process used, the outcome
expected from the original design (based on RFC~792, RFC~1122 and
RFC~1191), the outcome actually observed in a live network-namespace
lab, an analysis of any discrepancy between the two, and the assumptions
the attack or defense implicitly depends on. Several outcomes differed
materially from the initial design assumptions; these differences, and
their causes, are the primary findings of this report.
\end{abstract}

\tableofcontents
\newpage

\section{Introduction}

Both attacks exploit the same structural property of ICMP error
handling: an ICMP error message embeds a copy of the original packet's
IP and TCP headers (RFC~792) so the receiving host can match the error
to one of its own connections. Because this embedded header is
attacker-controlled once source-address spoofing is possible, a forged
ICMP message that mimics a real connection's headers can influence a TCP
stack's behavior without the attacker ever being on-path for that
connection.

\begin{itemize}
    \item \textbf{Attack 1} spoofs \textbf{ICMP Type 3, Code 3} (Port
    Unreachable) to trigger an RFC~1122 \emph{hard error}, which a
    vulnerable stack uses to abort an established connection immediately.
    \item \textbf{Attack 2} spoofs \textbf{ICMP Type 3, Code 4}
    (Fragmentation Needed, DF Set) with a small advertised Next-Hop MTU,
    to shrink a bulk sender's Path MTU (and therefore its MSS) and
    degrade throughput.
\end{itemize}

Both were implemented and tested in a controlled lab of three Linux
network namespaces (\texttt{client} 10.0.0.1, \texttt{server} 10.0.0.2,
\texttt{attacker} 10.0.0.3) joined by a software bridge, as specified in
the design report and built by \texttt{setup\_topology.sh}.

\section{Attack 1: ICMP Blind Connection-Reset}

\subsection{Attack Process}

\begin{enumerate}
    \item A client and server hold an ordinary, idle TCP connection open
    (via \texttt{ncat}).
    \item The attacker passively sniffs the bridge (via traffic
    mirroring configured in \texttt{setup\_topology.sh}) to observe the
    server's own outgoing segment and recover the client's ephemeral
    port and a sequence number known to be inside the server's current
    send window (\texttt{attacks/build\_packets.py::capture\_live\_state()}).
    \item The attacker crafts a spoofed ICMP Type 3/Code 3 packet
    (\texttt{attacks/build\_packets.py::build\_reset\_packet()}) and
    injects it with Scapy's \texttt{send()}
    (\texttt{attacks/inject\_reset.py}).
    \item The server's kernel receives the ICMP, matches it against the
    live socket via the embedded header, and --- if it treats the error
    as fatal --- aborts the connection.
\end{enumerate}

\paragraph{Packet structure used (corrected).} The design report's
original field table specified the embedded header as
\texttt{src=Client, dst=Server}. Empirically, this direction never
matched any socket on the server, because ICMP errors are delivered to
the \emph{original sender} of the embedded packet -- the server has no
socket where it is the destination of its own traffic. The corrected,
working structure is:

\begin{center}
\begin{tabular}{@{}lll@{}}
\toprule
\textbf{Layer} & \textbf{Field} & \textbf{Value} \\
\midrule
Outer IP & Source & 10.0.0.2 (spoofed -- impersonates the server) \\
Outer IP & Destination & 10.0.0.2 (the real server; the target) \\
ICMP & Type / Code & 3 / 3 (Port Unreachable) \\
Embedded IP & Source $\to$ Destination & 10.0.0.2 $\to$ 10.0.0.1 (server $\to$ client) \\
Embedded TCP & Source port $\to$ Dest.\ port & server's port $\to$ client's ephemeral port \\
Embedded TCP & Sequence number & sniffed from the server's own live outgoing segment \\
\bottomrule
\end{tabular}
\end{center}

\subsection{Expected Outcome}

Per the design report (Section 5.1), a successful attack should produce:
\begin{enumerate}
    \item The spoofed ICMP visible in the server's own \texttt{tcpdump} capture.
    \item The embedded sequence number falling inside the server's receive window.
    \item The server's TCP stack aborting the connection (RFC~1122 hard-error behavior).
    \item \texttt{ss -tn} on the server showing the socket gone.
    \item The client-side \texttt{ncat} process reporting a broken pipe / reset.
\end{enumerate}

\subsection{Actual Outcome}

Items 1 and 2 were confirmed on every run: the spoofed packet reliably
arrives and reliably passes the in-window check. \textbf{Items 3--5 did
not occur against a realistic target.} Against a connection held with
plain \texttt{ncat}, the same, correctly-formed packet was delivered but
the connection \textbf{survived} -- it remained \texttt{ESTAB} in
\texttt{ss -tn} with no observable effect.

\begin{center}
\begin{tabular}{@{}lll@{}}
\toprule
\textbf{Scenario} & \textbf{ICMP delivered / in-window?} & \textbf{Result} \\
\midrule
Plain \texttt{ncat} (realistic target) & Yes & \textbf{Survived} (still \texttt{ESTAB}) \\
\texttt{recverr\_server.py} (\texttt{IP\_RECVERR} set) & Yes & \textbf{Reset} (socket left \texttt{ESTAB}, typically \texttt{FIN-WAIT-2}) \\
\bottomrule
\end{tabular}
\end{center}

To rule out a defect in the injected packet itself, a second target,
\texttt{experiments/recverr\_server.py}, was built: an otherwise-identical
TCP server that explicitly sets the \texttt{IP\_RECVERR} socket option.
Against this target the identical attack \textbf{did} reliably reset the
connection (\texttt{ConnectionRefusedError} raised on the server, socket
observed leaving \texttt{ESTAB}), confirming the packet's direction,
sequence number, and checksums were all correctly formed.

\subsection{Discrepancy Analysis: Why the Outcome Differed}

Reading the Linux kernel source (\texttt{tcp\_v4\_err()} in
\texttt{net/ipv4/tcp\_ipv4.c}) explains the discrepancy. For a socket
that is already in the \texttt{ESTABLISHED} state, Linux only treats
this ICMP as an immediately-fatal RFC~1122 hard error if the owning
application has explicitly opted in via the \texttt{IP\_RECVERR} socket
option; otherwise the kernel silently downgrades it to a non-fatal
\texttt{sk\_err\_soft} and the connection is left untouched. Ordinary
tools such as \texttt{ncat} do not set this option, so they are immune
to this attack in its simplest form, while an application that does set
it (a legitimate, documented use case for applications that want richer
error reporting) remains vulnerable exactly as the design predicted.

This is a real, previously undocumented (in the design report) kernel
defense, additional to the two the design report already lists in
Section 6, and is reported as such in Section~5 of this report.

\subsection{Assumptions}

\begin{itemize}
    \item The attacker can observe the live connection's traffic (via
    mirroring in this lab; in a real network, via some other on-path
    vantage point, a shared broadcast segment, or coarse
    inference/guessing of the sequence number).
    \item Ingress/egress source-address filtering is \emph{not} enforced
    anywhere between the attacker and the target (\texttt{rp\_filter} is
    explicitly disabled in the lab to permit this).
    \item The targeted application either does not rely on the specific
    \texttt{IP\_RECVERR} behavior described above, or the report
    explicitly evaluates the case where it does (the positive control).
    \item The embedded sequence number remains inside the receiver's
    current window at the moment the forged packet is processed --
    trivially true here because the target connection is deliberately
    idle, so the window never moves.
\end{itemize}

\section{Attack 2: ICMP Blind PMTU Throughput Reduction}

\subsection{Attack Process}

\begin{enumerate}
    \item A client sends a sustained bulk transfer to the server via
    \texttt{iperf3}.
    \item The attacker sniffs the client's own outgoing segments to
    recover a live sequence number
    (\texttt{attacks/build\_packets.py::capture\_live\_state()}).
    \item The attacker crafts a spoofed ICMP Type 3/Code 4 packet with a
    small advertised Next-Hop MTU (576 bytes) and \emph{sprays} 300
    candidate packets, spaced 50{,}000 bytes apart in sequence space
    (\texttt{attacks/inject\_pmtu.py}), repeating this every few seconds
    for the duration of the attack.
    \item The client's kernel accepts the ICMP, updates its Path-MTU
    cache for the server, and shrinks its MSS accordingly.
\end{enumerate}

\paragraph{Packet structure used.} Unlike Attack 1, the design report's
direction for this attack was already correct, because the client is
both the connection's sender \emph{and} the intended target (MSS is a
per-sender cache):

\begin{center}
\begin{tabular}{@{}lll@{}}
\toprule
\textbf{Layer} & \textbf{Field} & \textbf{Value} \\
\midrule
Outer IP & Source & 10.0.0.1 (spoofed) \\
Outer IP & Destination & 10.0.0.1 (the real client; the target, as sender) \\
ICMP & Type / Code / Next-Hop MTU & 3 / 4 / 576 \\
Embedded IP & Source $\to$ Destination & 10.0.0.1 $\to$ 10.0.0.2 (client $\to$ server) \\
Embedded TCP & Source port $\to$ Dest.\ port & client's ephemeral port $\to$ 5201 \\
Embedded TCP & Sequence number & sniffed from the client's own live outgoing traffic \\
\bottomrule
\end{tabular}
\end{center}

\subsection{Expected Outcome}

Per the design report (Section 5.2):
\begin{enumerate}
    \item The spoofed ICMP visible in the client's capture.
    \item \texttt{ip route get} showing \texttt{mtu 576} (PMTU cache updated).
    \item TCP segment sizes dropping from $\sim$1460 bytes to $\le 576 - 40 = 536$ bytes.
    \item A measurable decrease in \texttt{iperf3} throughput.
\end{enumerate}

\subsection{Actual Outcome}

All four criteria were confirmed, across multiple independent runs.
Representative measured values:

\begin{center}
\begin{longtable}{@{}lll@{}}
\toprule
\textbf{Metric} & \textbf{Before injection} & \textbf{After injection} \\
\midrule
\endhead
\texttt{ss -tin} (\texttt{mss}/\texttt{pmtu}) & \texttt{mss:1448 pmtu:1500} & \texttt{mss:524 pmtu:576} \\
\texttt{ip route get 10.0.0.2} & no cached MTU & \texttt{cache expires <N>sec mtu 576} \\
Avg.\ TCP segment size & 1448 bytes & 531--542 bytes \\
\texttt{iperf3} throughput (run A) & 95.21 Mbps & 89.03 Mbps ($-6.5\%$) \\
\texttt{iperf3} throughput (run B) & 92.83 Mbps & 87.57 Mbps ($-5.7\%$) \\
ICMP packets observed & --- & exactly (spray rounds $\times$ 300) \\
\bottomrule
\end{longtable}
\end{center}

The ICMP-packet count matching \emph{exactly} (spray rounds $\times$
300) across every run is an internal consistency check confirming every
sprayed candidate was actually delivered, not merely that some
non-deterministic subset landed.

\subsection{Discrepancy Analysis: Why the Outcome Differed (Partially)}

The four qualitative criteria all held, but two aspects needed
correction from a naive first implementation, both already anticipated
in the design report's phrasing but underestimated in magnitude:

\begin{itemize}
    \item \textbf{A single sequence-number guess reliably fails.}
    Unlike Attack 1's idle target, this connection is an active bulk
    transfer, so the kernel's in-window check
    (\texttt{between(seq, snd\_una, snd\_nxt)}) -- which applies to
    \emph{every} ICMP type/code, not only the reset case -- must be
    satisfied against a window that is only \texttt{cwnd*mss} wide
    (measured at $\sim$100--150\,KB) and sliding forward at the full
    transfer rate. By the time one sniffed value passes through
    \texttt{tcpdump}, Python, and Scapy's \texttt{send()}, it is almost
    always already acknowledged and out of window (confirmed via the
    \texttt{TcpExt:OutOfWindowIcmps} counter in \texttt{/proc/net/netstat}).
    Spraying 300 candidates across $\sim$15\,MB of sequence space, in
    steps smaller than the in-flight window, is what reliably lands at
    least one packet inside the live window. This matches the design
    report's own (plural) wording, ``sampled seq numbers,'' and the
    technique used in Watson's original blind TCP injection research.
    \item \textbf{The throughput drop is real but modest ($5$--$7\%$),
    not dramatic}, because the lab's client--server link is
    byte-rate-limited to 100\,Mbit (\texttt{tc tbf}) for realism. At a
    fixed byte rate, a smaller MSS mainly shifts more of every
    rate-capped packet toward the fixed 40-byte IP+TCP header instead of
    payload ($\sim$97\% payload efficiency at MSS 1448 vs.\ $\sim$93\%
    at MSS 536) -- a real, explainable effect of the lab's shaping
    choice, not evidence the attack is weak.
\end{itemize}

\subsection{Assumptions}

\begin{itemize}
    \item The attacker can observe the live connection's own outgoing
    traffic to recover a plausible base sequence number.
    \item No ingress/egress source-address filtering exists between the
    attacker and the target.
    \item The attacker can send the spray fast enough, and repeat it
    often enough, to keep at least one candidate inside the constantly
    shifting in-window range -- a persistent-attacker assumption,
    explicitly acknowledged in \S 4 below when discussing why a
    self-healing defense can be undermined by re-spraying.
    \item The target's application does not itself detect or route
    around the reduced MSS in a way that this report has not modeled
    (e.g.\ redundant multipath).
\end{itemize}

\section{Defense Mechanisms}

Two defenses named in the design report (Section 6) were implemented
and evaluated, plus one already-listed defense that was confirmed as an
emergent property of the attacks themselves rather than a separately
built mechanism.

\subsection{Sequence-Number Validation (Design Report Defense 1)}

\paragraph{Mechanism.} Modern Linux kernels reject any ICMP error whose
embedded sequence number falls outside the target connection's current
window, regardless of whether the source address is genuine.

\paragraph{Expected outcome.} Blind sequence-number guessing becomes
substantially harder against an active connection.

\paragraph{Actual outcome.} \worked. This defense was not implemented
separately -- it is already active, by default, in the lab's stock
Linux kernel, and its effect is precisely what forced Attack 2's design
away from a single guess toward the 300-candidate spray described in
Section~3.3. Its presence is confirmed indirectly but unambiguously by
the \texttt{TcpExt:OutOfWindowIcmps} counter incrementing on single-shot
attempts.

\subsection{Ingress Filtering / BCP~38 (Design Report Defense 2)}

\paragraph{Mechanism.} A border router or switch port refuses to
forward any packet whose source address does not belong to the network
segment it physically arrived from. Simulated in this lab with a
\texttt{tc} filter on the bridge port the attacker's traffic physically
arrives on (\texttt{veth-at-br}), permitting only
\texttt{src=10.0.0.3} (the attacker's genuine address) and dropping
everything else:

\begin{verbatim}
tc qdisc add dev veth-at-br handle ffff: ingress
tc filter add dev veth-at-br parent ffff: protocol ip prio 1 u32 \
    match ip src 10.0.0.3/32 action ok
tc filter add dev veth-at-br parent ffff: protocol ip prio 2 u32 \
    match u32 0 0 action drop
\end{verbatim}

\paragraph{Expected outcome.} Both attacks should be blocked completely,
regardless of the target's own TCP-stack-specific behavior, since the
spoofed packet is stopped before it ever reaches the bridge.

\paragraph{Actual outcome.} \worked. Re-running both attacks with the
filter active produced no ICMP at either target and no change to
connection state, MSS, PMTU, segment size, or throughput in either
attack. Critically, the \texttt{recverr\_server.py} scenario -- which
\emph{always} resets without this defense -- \textbf{survived} with the
filter active, isolating that the defense operates at the network layer
and is independent of the target application's socket options. This
was the single most conclusive and reliable defense evaluated.

\begin{center}
\begin{tabular}{@{}lll@{}}
\toprule
\textbf{Scenario (filter active)} & \textbf{ICMP delivered?} & \textbf{Result} \\
\midrule
Attack 1, plain \texttt{ncat} & No & Survived (unchanged) \\
Attack 1, \texttt{recverr\_server.py} & No & \textbf{Survived} (previously reset without the filter) \\
Attack 2, \texttt{iperf3} bulk transfer & No & \texttt{mss}/\texttt{pmtu} unchanged; throughput unaffected \\
\bottomrule
\end{tabular}
\end{center}

\subsection{PMTU Cache Aging / PLPMTUD (Design Report's Implementation Plan, P7)}

\paragraph{Mechanism (as specified).} The design report's Implementation
Plan instructs: ``Enable \texttt{tcp\_mtu\_probing=2}. Repeat P4 and P5.
Record whether the kernel rejects the forged ICMP messages.'' The
intent is that Linux's Packetization-Layer PMTU Discovery (PLPMTUD)
would cause the client to re-verify its true path MTU and recover from
the forged value on its own.

\paragraph{Expected outcome.} A single spoofed PMTU packet degrades the
connection's MSS momentarily, then the client's own kernel-level
recovery restores \texttt{mss}/\texttt{pmtu} to their original values
within a short window, without any further attacker or network-layer
intervention.

\paragraph{Actual outcome.} \failed. Two successive, independently
testable hypotheses about \emph{how} this recovery should occur were
each implemented and tested, and \textbf{neither produced the expected
recovery} within the tested window:

\begin{enumerate}
    \item \textbf{First hypothesis:} \texttt{net.ipv4.tcp\_mtu\_probing=2}
    together with a shortened \texttt{net.ipv4.tcp\_probe\_interval}
    (5\,s) would cause active re-probing and recovery. \emph{Result:}
    \texttt{mss}/\texttt{pmtu} remained pinned at the forged
    \texttt{524}/\texttt{576} for over 20 continuous seconds. On
    inspection, \texttt{tcp\_mtu\_probing} governs detection of ICMP
    \emph{black holes} (a router that silently drops oversized packets
    and never replies at all) and has no role in re-verifying a PMTU
    value the kernel already accepted from an ICMP message it did
    receive.
    \item \textbf{Second hypothesis:} the route-cache PMTU
    \emph{exception} created when the ICMP lands (visible as
    \texttt{cache expires <N>sec mtu 576} in \texttt{ip route get}) ages
    out via \texttt{net.ipv4.route.mtu\_expires} (600\,s by default);
    this was shortened to 8\,s to make recovery observable within a
    reasonable test window. \emph{Result:} even after confirming the
    exception was created with the shortened timer, \texttt{mss}/\texttt{pmtu}
    were still pinned at \texttt{524}/\texttt{576} at every checkpoint
    out to $t{+}40$\,s.
\end{enumerate}

\paragraph{Status at the time of writing.} This defense is
\open. A diagnostic step was added (checking whether the route-table
exception itself has actually been evicted by $t{+}40$\,s, then forcibly
flushing the route cache to determine whether an already-established,
continuously busy socket needs an explicit external trigger to
re-resolve its cached PMTU) but had not yet produced a confirmed root
cause at submission time. The most plausible working explanation is
that Linux caches a socket's last-known PMTU
(\texttt{icsk\_pmtu\_cookie}) and only re-validates it opportunistically
--- not automatically upon a route-table exception's expiry elsewhere
--- so a continuously active socket may never receive that trigger on
its own.

\subsection{Summary: Design-Report Defenses vs.\ Evaluated Reality}

\begin{center}
\begin{longtable}{@{}p{4.2cm}p{2.6cm}p{7.5cm}@{}}
\toprule
\textbf{Defense} & \textbf{Verdict} & \textbf{Notes} \\
\midrule
\endhead
Sequence-number validation & \worked & Inherent to the stock kernel; confirmed indirectly via the need for Attack 2's 300-packet spray. \\
Ingress filtering (BCP~38) & \worked & Blocks both attacks completely; independent of target application behavior; the single most reliable defense tested. \\
\texttt{tcp\_mtu\_probing=2} (P7, as literally specified) & \failed & Governs ICMP black-hole detection, not recovery from an \emph{accepted} forged ICMP; produced no observed recovery. \\
PMTU route-cache aging (follow-up hypothesis, not in the original design report) & \open & Shortened \texttt{mtu\_expires} alone also failed to produce recovery within 40\,s; root cause still being isolated via a route-cache-flush diagnostic. \\
\texttt{IP\_RECVERR} gating (discovered during Attack 1 testing, not in the original design report) & \worked & Not a designed defense, but a real, already-built-in Linux behavior that defeats Attack 1 against any application that does not explicitly request richer ICMP error delivery. \\
\bottomrule
\end{longtable}
\end{center}

\section{Consolidated Assumptions}

The following assumptions are shared across both attacks and are
necessary preconditions for either to succeed at all:

\begin{enumerate}
    \item \textbf{No source-address filtering} exists anywhere between
    the attacker and the target -- disabled explicitly in this lab
    (\texttt{rp\_filter=0}) to permit spoofing, and shown in
    Section~4.2 to fully neutralize both attacks when reintroduced.
    \item \textbf{The attacker can observe enough of the live
    connection} (source/destination ports and an in-window sequence
    number) to construct a convincing embedded header -- provided here
    by passive traffic mirroring, but achievable in principle by other
    means (e.g.\ a shared broadcast segment, or brute-force sequence
    guessing at a cost proportional to the window size).
    \item \textbf{The target's TCP stack processes the embedded
    header using the packet's actual direction semantics} -- i.e.\ ICMP
    errors are delivered to, and matched against, the original sender
    of the embedded packet. This governed which side (server vs.\
    client) each attack had to impersonate, and was the single
    correction required to the design report's original packet tables.
    \item \textbf{For Attack 1 specifically}: the target application
    either sets \texttt{IP\_RECVERR} (uncommon, e.g.\ specialized error
    reporting) or the report separately evaluates that narrower case as
    a positive control, since the attack does not succeed against a
    typical established connection without it.
    \item \textbf{For Attack 2 specifically}: the attacker must be able
    to repeat the spray persistently against an undefended target,
    since (regardless of which recovery mechanism eventually turns out
    to govern it) a defended or self-healing target can undo a
    single-shot attack given enough time.
\end{enumerate}

\section{Conclusion}

Both attacks were successfully implemented and demonstrated the core
mechanism described in the design report: a spoofed ICMP error,
carrying a forged embedded header, can influence a real TCP connection's
behavior without the attacker being on-path. Attack 2 (PMTU reduction)
matched its design-time expectations closely once corrected for the
in-window sequence constraint. Attack 1 (connection reset) required a
larger correction: it does \emph{not} reliably work against a realistic
target, and its true precondition (\texttt{IP\_RECVERR}) is narrower
than the design report assumed -- a materially different, and arguably
more interesting, result than a simple confirmation would have been.

Of the two evaluated defenses, ingress filtering proved to be a
complete and reliable countermeasure against both attacks, while the
kernel self-healing mechanism the design report's Implementation Plan
specified (\texttt{tcp\_mtu\_probing=2}) did not produce its expected
effect under either of two tested hypotheses for why it should.
Reporting this negative result, rather than only reporting successes,
is treated in this report as itself a legitimate and important finding:
not every plausible, textbook-named defense behaves as its name
suggests once it is actually measured.

\begin{thebibliography}{9}
\bibitem{rfc792} J. Postel, \emph{Internet Control Message Protocol}, RFC~792, Sept.\ 1981.
\bibitem{rfc1122} R. Braden (Ed.), \emph{Requirements for Internet Hosts --- Communication Layers}, RFC~1122, Oct.\ 1989.
\bibitem{rfc1191} J. Mogul and S. Deering, \emph{Path MTU Discovery}, RFC~1191, Nov.\ 1990.
\bibitem{watson2004} P. Watson, \emph{Slipping in the Window: TCP Reset Attacks}, CanSecWest, 2004.
\bibitem{rfc5961} A. Ramaiah, R. Stewart, M. Dalal, \emph{Improving TCP's Robustness to Blind In-Window Attacks}, RFC~5961, Aug.\ 2010.
\bibitem{rfc2827} P. Ferguson and D. Senie, \emph{Network Ingress Filtering: Defeating Denial of Service Attacks which employ IP Source Address Spoofing}, BCP~38 / RFC~2827, May 2000.
\bibitem{rfc6633} F. Gont, \emph{Deprecation of ICMP Source Quench Messages}, RFC~6633, May 2012.
\end{thebibliography}

\end{document}
